\documentclass[11pt,a4paper]{article}

\usepackage[T1]{fontenc}
\usepackage[utf8]{inputenc}
\usepackage{mathptmx}
\usepackage{amsmath,amssymb,amsthm}
\usepackage{geometry}
\usepackage{microtype}
\usepackage{hyperref}
\usepackage{booktabs}
\usepackage{enumitem}
\usepackage{setspace}
\usepackage{titlesec}
\usepackage{abstract}
\usepackage{natbib}
\usepackage{xcolor}

\geometry{top=1.1in, bottom=1.1in, left=1.15in, right=1.15in}

\titleformat{\section}{\normalfont\bfseries\large}{\thesection}{1em}{}
\titleformat{\subsection}{\normalfont\bfseries\normalsize}{\thesubsection}{1em}{}
\titleformat{\subsubsection}{\normalfont\bfseries\itshape\normalsize}{\thesubsubsection}{1em}{}

\newtheorem{theorem}{Theorem}[section]
\newtheorem{lemma}[theorem]{Lemma}
\newtheorem{proposition}[theorem]{Proposition}
\newtheorem{corollary}[theorem]{Corollary}
\newtheorem{definition}[theorem]{Definition}
\newtheorem{remark}[theorem]{Remark}

\newcommand{\Pact}{\mathcal{P}_{\mathrm{act}}}
\newcommand{\SBDG}{S_{\mathrm{BDG}}}
\newcommand{\LLC}{\mathrm{LLC}}
\newcommand{\dS}{\Delta S}

\hypersetup{
  colorlinks=true, linkcolor=blue, citecolor=blue, urlcolor=blue,
  pdftitle={Relational Actualism: The Einstein Field Equations as the Unique Macroscopic Limit of the Actualization Graph (RACL) -- doi:10.5281/zenodo.19270020},
  pdfauthor={Joshua F. Sandeman}
}

\begin{document}

% ── Title block ───────────────────────────────────────────────────────────────
\begin{center}
  {\LARGE\bfseries Relational Actualism: The Einstein Field Equations\\
  as the Unique Macroscopic Limit of the Actualization Graph\par}
  \vspace{6pt}
  {\large\bfseries Local Ledger Conservation, P$_{\mathrm{act}}$ Preservation,\\
  and the Lovelock Derivation of GR with $\Lambda = 0$\par}
  \vspace{14pt}
  {\large Joshua F.\ Sandeman$^{1*}$\par}
  \vspace{4pt}
  {\normalsize $^{1}$Independent Researcher, Salem, OR, USA.\par}
  {\normalsize $^{*}$sansuikyo@gmail.com\par}
  \vspace{8pt}
  {\small\textit{Preprint available at}~\href{https://doi.org/10.5281/zenodo.19270020}{doi.org/10.5281/zenodo.19270020}\par}
  \vspace{14pt}
\end{center}

% ── Abstract ─────────────────────────────────────────────────────────────────
\begin{abstract}
We prove that the Einstein field equations with $\Lambda = 0$ are the
\emph{unique} consistent macroscopic description of the Relational
Actualism causal graph (RA-CSG), with no free parameters.
The derivation combines four ingredients: the Lean~4-verified
Local Ledger Condition (LLC); the Benincasa-Dowker theorem
$\langle\SBDG\rangle \to (l_P^2/4)R$; Lovelock's uniqueness theorem;
and vacuum energy suppression via the actualization projector $\Pact$,
which forces $\Lambda = 0$ structurally. The Bianchi identity holds
by geometry.

Two results are new. First, a \emph{BDG locality lemma}: the BDG
score depends only on ancestors within proper-time window
$\tau_{\mathrm{BDG}}(\lambda) = 4\,\Gamma(5/4)\,(12/\pi^2)^{1/4}\,
\lambda^{-1/4} \to 0$ as $\lambda \to \infty$, confirming locality
of the continuum limit. Second, the \emph{$\Pact$ conservation
theorem}: $\nabla_\mu(\Pact T^{\mu\nu}) = 0$, proved via Leibniz
decomposition using the LLC and the Lean-verified frame-independence
of the relative entropy $\Delta S(\rho\|\sigma_0)$.

The proof is RA-native: the LLC---a vertex-level charge balance
condition verified in Lean~4 with zero \texttt{sorry} tags---forces
the macroscopic field equation to be the Einstein equation.
GR is an output of RA, not an input.

Where RA departs from GR is where $\Pact T^{\mu\nu}$ departs from
$T^{\mu\nu}$: in sparse regions where the on-shell fraction
$\lambda_{\rm local}/\lambda_{\rm ref} < 1$. These departures are
classical actualization-density gradient effects, not
Planck-scale corrections, and correspond to the Hubble tension
and galactic rotation curves.
\end{abstract}

\noindent\textbf{Keywords:} causal sets; Einstein field equations; Local Ledger Condition; Lovelock uniqueness theorem; Relational Actualism; formal verification; Lean 4; quantum gravity; emergent spacetime; cosmological constant

\medskip
\noindent\textbf{Ethics statement.}
This article does not contain any studies with human participants
or animals performed by the author.



% ── 1. Introduction ──────────────────────────────────────────────────────────
\section{Introduction}
\label{sec:intro}

A foundational question for any discrete approach to quantum gravity
is: why does general relativity emerge at macroscopic scales? The
standard answer in causal set theory invokes the Benincasa-Dowker
(BDG) action~\citep{Benincasa2010}, which converges to the
Einstein-Hilbert action in the dense-graph limit. But convergence of
an action is not the same as uniqueness of the resulting field
equation, and neither result directly answers the deeper question of
\emph{why} the macroscopic field equation has no free parameters.

Relational Actualism (RA)~\citep{Sandeman2026RAQM} provides a
different starting point. RA grounds physics in irreversible on-shell
actualization events, each writing a permanent directed edge into a
growing causal Directed Acyclic Graph (DAG). The primitive is not an
action but a \emph{local balance condition}: the Local Ledger
Condition (LLC), which requires that the sum of incoming charges
equals the sum of outgoing charges at every actualization vertex.
The LLC is machine-verified in Lean~4 with zero \texttt{sorry}
tags~\citep{Sandeman2026RAGC}.

This paper proves that given the LLC alone --- together with three
published theorems --- the Einstein field equations are the unique
macroscopic field equation of the RA-CSG, with $\Lambda = 0$ and no
remaining free parameters. The argument does not treat GR as a target
to be recovered from a limiting procedure. It shows that the LLC
\emph{forces} the macroscopic dynamics to be GR in the regime where
the actualization-density projector $\Pact$ acts as the identity on
$T^{\mu\nu}$. Where $\Pact \neq \mathrm{id}$, the field equations
predict specific, observable departures from GR that match the Hubble
tension and galactic rotation curves~\citep{Sandeman2026RAGC,
Sandeman2026RADM}.

\medskip\noindent
\textbf{Structure.}
Section~\ref{sec:setup} establishes notation and the three imported
theorems. Section~\ref{sec:locality} proves the BDG locality lemma.
Section~\ref{sec:P_act} proves the $\Pact$ conservation theorem.
Section~\ref{sec:main} assembles the main derivation.
Section~\ref{sec:departures} characterises the RA departures from GR.
Section~\ref{sec:lean} discusses the Lean~4 verification status.
Section~\ref{sec:related} situates the work in the existing literature.

% ── 2. Setup ─────────────────────────────────────────────────────────────────
\section{Setup and Imported Theorems}
\label{sec:setup}

\subsection{The RA-CSG and the LLC}

Let $G = (V, E)$ be the RA causal graph: a growing directed acyclic
graph in which each vertex $v \in V$ is an on-shell actualization
event and each directed edge $(u,v) \in E$ records a causal
relationship $u \prec v$. Each vertex carries a charge vector
$\mathbf{q}(v) \in \mathbb{R}^n$ encoding conserved quantities
(energy-momentum, electric charge, baryon number, etc.).

\begin{definition}[Local Ledger Condition]
\label{def:LLC}
For every vertex $v \in V$:
\begin{equation}
  \sum_{\substack{u \in V \\ (u,v)\in E}} \mathbf{q}(u)
  \;=\;
  \sum_{\substack{w \in V \\ (v,w)\in E}} \mathbf{q}(w).
\end{equation}
\end{definition}

\noindent
The LLC is Lean~4-verified with zero \texttt{sorry} tags in
\texttt{RA\_Proofs\_Lean4.lean}. The Graph Cut
Theorem~\citep{Sandeman2026RAGC} shows the LLC holds independently on
each side of a causal severance, with no \texttt{sorry} in the Lean
proof.

\subsection{The actualization projector}

The actualization criterion is $\dS(\rho\|\sigma_0) > 0$, where
$\rho$ is the physical state and $\sigma_0$ the vacuum~\citep{Sandeman2026RAQM}.
Define:
\begin{equation}
  \Pact = \Theta\!\bigl(\dS(\rho\|\sigma_0)\bigr),
\end{equation}
where $\Theta$ is the Heaviside step function. The operator $\Pact$
projects onto on-shell (actualized) states and annihilates the vacuum.
Its frame-independence is Lean~4-proved:
\texttt{frame\_independence} in \texttt{RA\_Proofs\_Lean4.lean},
using the continuous functional calculus unitary-conjugation lemma
from Lean-QuantumInfo~\citep{Meiburg2025LQI}.

The vacuum energy suppression theorem~\citep{Sandeman2026RAGC} gives:
\begin{equation}
  \bigl\langle 0 \bigm| \Pact\, T^{\mu\nu} \bigm| 0 \bigr\rangle = 0.
  \label{eq:vacuum_suppression}
\end{equation}
This follows directly from $\Pact|0\rangle = 0$: the vacuum contains
no actualized states. It is a \emph{structural} result, not
fine-tuning.

\subsection{Three imported theorems}

The derivation imports three published results without reproving them.

\begin{theorem}[Rideout-Sorkin density limit, \citep{RideoutSorkin2000}]
\label{thm:rs}
In the dense-graph limit $\lambda \to \infty$ of the Poisson
Classical Sequential Growth (CSG) dynamics, the mean of any local
observable on the causal set converges to a well-defined continuum
limit.
\end{theorem}

\begin{theorem}[Benincasa-Dowker, \citep{Benincasa2010}]
\label{thm:bd}
For a Poisson distribution of vertices at density $\rho$ sprinkled
into a $d{=}4$ Lorentzian manifold with Ricci scalar $R(x)$:
\begin{equation}
  \bigl\langle \SBDG(v) \bigr\rangle
  \;=\; \frac{l_P^2}{4}\,R\bigl(x(v)\bigr)
        \;+\; \mathcal{O}(\rho^{-1/2}).
\end{equation}
The correction vanishes as $\rho \to \infty$.
\end{theorem}

\begin{theorem}[Lovelock, \citep{Lovelock1971}]
\label{thm:lovelock}
In four spacetime dimensions, the unique symmetric
divergence-free tensor of rank $(0,2)$ that is local, covariant,
and built from the metric and its derivatives up to second order is:
\begin{equation}
  H_{\mu\nu} = \alpha\, G_{\mu\nu} + \Lambda\, g_{\mu\nu},
\end{equation}
where $G_{\mu\nu}$ is the Einstein tensor and $\alpha$, $\Lambda$
are constants.
\end{theorem}

% ── 3. BDG Locality Lemma ────────────────────────────────────────────────────
\section{BDG Locality Lemma}
\label{sec:locality}

The BDG score at vertex $v$ involves $N_k(v)$ for $k = 1,2,3,4$,
where $N_k(v)$ counts $k$-element causal intervals in
$\mathrm{past}(v)$. A $k$-element interval is a chain
$u_0 \prec u_1 \prec \cdots \prec u_{k-1} \prec v$ of length $k$.

\begin{lemma}[BDG 4-hop locality]
\label{lem:locality}
At Poisson vertex density $\lambda$ (per unit 4-volume in Planck
units), the BDG score $\SBDG(v)$ depends only on vertices within a
proper-time window
\begin{equation}
  \tau_{\mathrm{BDG}}(\lambda)
  \;=\; 4\,\Gamma\!\left(\tfrac{5}{4}\right)
        \left(\frac{12}{\pi^2}\right)^{\!1/4}
        \lambda^{-1/4}.
  \label{eq:tau_BDG}
\end{equation}
As $\lambda \to \infty$: $\tau_{\mathrm{BDG}}(\lambda) \to 0$.
\end{lemma}

\begin{proof}
The 4-volume of a causal diamond of proper time $\tau$ in 4D
Minkowski space is $V_4(\tau) = (\pi^2/12)\,\tau^4$. In a Poisson
causal set, the probability that a causal diamond of proper time $\tau$
contains no intermediate vertex is $e^{-\lambda V_4(\tau)}$. The
expected proper time of a single causal hop (direct causal relation
with no intermediate element) is therefore:
\begin{equation}
  \langle\tau_{\mathrm{hop}}\rangle
  \;=\; \int_0^\infty \tau\,
        \frac{d}{d\tau}\!\left[1 - e^{-\lambda V_4(\tau)}\right]
        e^{-\lambda V_4(\tau)}\, d\tau
  \;\sim\; \Gamma\!\left(\tfrac{5}{4}\right)
           \left(\frac{12}{\pi^2\lambda}\right)^{\!1/4}.
\end{equation}
A $k$-hop interval spans at most $k$ independent hops, giving total
proper time at most $k\,\langle\tau_{\mathrm{hop}}\rangle$. For
$k \leq 4$:
\begin{equation}
  \tau_{\mathrm{BDG}}(\lambda)
  = 4\,\langle\tau_{\mathrm{hop}}\rangle
  = 4\,\Gamma\!\left(\tfrac{5}{4}\right)
    \left(\frac{12}{\pi^2}\right)^{\!1/4}
    \lambda^{-1/4}.
\end{equation}
Since $\tau_{\mathrm{BDG}}(\lambda) \to 0$ as $\lambda \to \infty$,
the BDG score becomes a function of fields within an arbitrarily small
neighbourhood of $x(v)$: a local function. \hfill$\square$

\medskip\noindent
\emph{Numerical verification.} The scaling $\tau_{\mathrm{hop}} \propto
\lambda^{-1/4}$ was confirmed computationally: for each tenfold
increase in $\lambda$, the ratio $\tau_{\mathrm{hop}}(10\lambda) /
\tau_{\mathrm{hop}}(\lambda) = 10^{-1/4} \approx 0.5623$, matching
the theoretical prediction to four significant figures across four
decades of density.
\end{proof}

\begin{remark}
Lemma~\ref{lem:locality} implies that in the dense-graph limit, the
BDG action is a \emph{local} functional of the metric, precisely as
required for the Benincasa-Dowker theorem (Theorem~\ref{thm:bd}) and
for the Lovelock argument (Theorem~\ref{thm:lovelock}) to apply.
\end{remark}

% ── 4. P_act Conservation Theorem ────────────────────────────────────────────
\section{P\textsubscript{act} Conservation Theorem}
\label{sec:P_act}

\begin{theorem}[$\Pact$ conservation]
\label{thm:P_act}
Given the LLC (Lean-verified), the frame-independence of $\dS(\rho\|\sigma_0)$
(Lean-verified), and the vacuum energy suppression theorem
(proved in \citep{Sandeman2026RAGC}), the actualization-projected
stress tensor satisfies:
\begin{equation}
  \nabla_\mu\!\bigl(\Pact\,T^{\mu\nu}\bigr) = 0.
  \label{eq:P_act_conservation}
\end{equation}
\end{theorem}

\begin{proof}
By the Leibniz rule for covariant derivatives:
\begin{equation}
  \nabla_\mu\!\bigl(\Pact\,T^{\mu\nu}\bigr)
  = \bigl(\nabla_\mu \Pact\bigr)\,T^{\mu\nu}
    + \Pact\,\bigl(\nabla_\mu T^{\mu\nu}\bigr).
  \label{eq:leibniz}
\end{equation}

\textbf{Second term.}
The LLC at every vertex gives
$\sum_{\mathrm{in}} \mathbf{q} = \sum_{\mathrm{out}} \mathbf{q}$.
By Theorem~\ref{thm:rs} (Rideout-Sorkin), the dense-graph limit of
this discrete balance condition is the continuous conservation equation
$\nabla_\mu T^{\mu\nu} = 0$. Therefore
$\Pact(\nabla_\mu T^{\mu\nu}) = \Pact \cdot 0 = 0$.

\textbf{First term.}
Since $\Pact = \Theta(\dS(\rho\|\sigma_0))$, the chain rule gives:
\begin{equation}
  \nabla_\mu \Pact
  = \delta\!\bigl(\dS(\rho\|\sigma_0)\bigr)
    \cdot \nabla_\mu\!\bigl(\dS(\rho\|\sigma_0)\bigr).
\end{equation}
The frame-independence of $\dS$ (Lean-proved) ensures that
$\nabla_\mu\dS$ is a well-defined covariant tensor. The distributional
derivative $\delta(\dS)$ has support precisely where $\dS = 0$.

By Klein's inequality, $\dS(\rho\|\sigma_0) = 0$ if and only if
$\rho = \sigma_0$ (the vacuum state). At $\rho = \sigma_0$:
$T^{\mu\nu} = \langle\sigma_0|T^{\mu\nu}|\sigma_0\rangle = 0$
by the vacuum energy suppression theorem~\eqref{eq:vacuum_suppression}.

The standard distributional identity applies: if $f$ is smooth and
$f(x_0) = 0$, then $f(x)\,\delta(x - x_0) = f(x_0)\,\delta(x-x_0) = 0$.
With $f = T^{\mu\nu}$ and $x_0 = \{\dS = 0\}$:
\begin{equation}
  \bigl(\nabla_\mu \Pact\bigr)\,T^{\mu\nu}
  = \delta(\dS)\cdot\nabla_\mu(\dS)\cdot T^{\mu\nu}
  = \delta(\dS)\cdot\nabla_\mu(\dS)\cdot 0 = 0.
\end{equation}

Substituting both terms into~\eqref{eq:leibniz}:
$\nabla_\mu(\Pact T^{\mu\nu}) = 0 + 0 = 0$. \hfill$\square$
\end{proof}

\begin{remark}
The frame-independence of $\dS$ (Lean-proved) is load-bearing: without
covariance of $\Pact$, the operator $\nabla_\mu\Pact$ would be
observer-dependent and the decomposition~\eqref{eq:leibniz} would not
yield a covariant result. The Lean proof of \texttt{frame\_independence}
in \texttt{RA\_Proofs\_Lean4.lean} is what makes this step rigorous.
\end{remark}

% ── 5. Main Derivation ───────────────────────────────────────────────────────
\section{Main Derivation: Uniqueness of the RA Field Equation}
\label{sec:main}

\begin{theorem}[RA field equation uniqueness]
\label{thm:main}
The Einstein field equations with $\Lambda = 0$,
\begin{equation}
  G_{\mu\nu} = 8\pi G\,\Pact\!\left[T_{\mu\nu}\right],
  \label{eq:RA_EFE}
\end{equation}
are the \emph{unique} consistent macroscopic field equation of the
RA-CSG. No free parameters remain.
\end{theorem}

\begin{proof}
The proof assembles the seven links in the chain
(A)$\to$(B)$\to$(C)$\to$(D)$\to$(E)$\to$(F)$\to$(G).

\medskip
\textbf{(A) LLC $\to$ conservation.}
By the LLC (Lean-verified) and Theorem~\ref{thm:rs}:
$\nabla_\mu(\Pact T^{\mu\nu}) = 0$ (Theorem~\ref{thm:P_act}).

\medskip
\textbf{(B) BDG action $\to$ local tensor.}
By Lemma~\ref{lem:locality}: $\tau_{\mathrm{BDG}} \to 0$ as
$\lambda \to \infty$, so $\SBDG(v)$ is a local function of the metric
at $x(v)$ in the dense-graph limit.

\medskip
\textbf{(C) Benincasa-Dowker.}
By Theorem~\ref{thm:bd}: the dense-graph mean
$\langle\SBDG\rangle = (l_P^2/4)R$, which is the Einstein-Hilbert
action density. Extremising $\int\langle\SBDG\rangle\,d^4x$ over
the metric yields $G_{\mu\nu} = 0$ in vacuum.

\medskip
\textbf{(D) Adding the matter source.}
Adding $\Pact T_{\mu\nu}$ as the matter source on the right:
\begin{equation}
  H_{\mu\nu} = \kappa\,\Pact T_{\mu\nu},
\end{equation}
where $H_{\mu\nu}$ is the left-hand side tensor to be determined.

\medskip
\textbf{(E) Lovelock uniqueness.}
The source $\Pact T_{\mu\nu}$ is divergence-free by step (A).
Therefore $H_{\mu\nu}$ must also be divergence-free:
$\nabla^\mu H_{\mu\nu} = \kappa\,\nabla^\mu(\Pact T_{\mu\nu}) = 0$.

By (B) and (C), $H_{\mu\nu}$ is local and second-order in the metric
(the BDG action is second-order by the Benincasa-Dowker theorem).

By Theorem~\ref{thm:lovelock} (Lovelock), the unique such tensor is:
$H_{\mu\nu} = \alpha\,G_{\mu\nu} + \Lambda\,g_{\mu\nu}$.

\medskip
\textbf{(F) Vacuum energy suppression.}
By~\eqref{eq:vacuum_suppression}: $\Pact|0\rangle = 0$, so the
vacuum contributes zero to the right-hand side. In standard GR, the
equivalent statement requires fine-tuning $\Lambda$ against the QFT
vacuum energy (the cosmological constant problem). In RA, $\Lambda = 0$
is structurally forced: the only states that source the metric are
actualized (on-shell) states, and the vacuum contains none.

With $\Lambda = 0$: $H_{\mu\nu} = \alpha\,G_{\mu\nu}$, giving:
\begin{equation}
  G_{\mu\nu} = \frac{\kappa}{\alpha}\,\Pact T_{\mu\nu}.
\end{equation}
Setting $\kappa/\alpha = 8\pi G$ recovers the Newtonian limit.

\medskip
\textbf{(G) Bianchi identity.}
$\nabla^\mu G_{\mu\nu} \equiv 0$ holds by the contracted Bianchi
identity, which is a geometric consequence of the Jacobi identity for
covariant derivatives on any Riemannian manifold. It holds for
\emph{any} $G_{\mu\nu}$ constructed from the Riemann tensor, and does
not depend on the field equation. No separate proof is required; it
is definitionally true.

\medskip
Combining (A)--(G): equation~\eqref{eq:RA_EFE} is the unique
consistent macroscopic field equation with no free parameters.
\hfill$\square$
\end{proof}

\begin{remark}[What is proved vs imported]
The derivation uses three imported published results:
Rideout-Sorkin~\citep{RideoutSorkin2000}, Benincasa-Dowker~\citep{Benincasa2010},
and Lovelock~\citep{Lovelock1971}. The RA-original contributions are:
the LLC (Lean-verified), the $\Pact$ conservation theorem (proved
above), the BDG locality lemma (proved above), and the vacuum
suppression theorem (proved in~\citep{Sandeman2026RAGC}).
The assembly of these ingredients into the uniqueness argument is new.
\end{remark}

% ── 6. Departures from GR ────────────────────────────────────────────────────
\section{Where RA Departs from GR}
\label{sec:departures}

Theorem~\ref{thm:main} establishes that~\eqref{eq:RA_EFE} is the
exact RA field equation. It reduces to standard GR when
$\Pact T_{\mu\nu} \approx T_{\mu\nu}$, i.e., when essentially all
of the matter is on-shell. This holds when the local actualization
density $\lambda_{\rm local}$ is high: in dense astrophysical
environments, nearly every mass-energy configuration is on-shell, and
GR is recovered to high precision.

When $\lambda_{\rm local}$ varies spatially, the departures from GR
are governed by the gradient of the ratio $\lambda_{\rm local}/\lambda_{\rm ref}$
(where $\lambda_{\rm ref}$ is the reference density in the dense
limit). At leading order in weak fields:
\begin{equation}
  \nabla^2\Phi_{\rm RA}
  = 4\pi G\bigl[\rho_A + \rho_\lambda\bigr],
  \quad
  \rho_\lambda = \frac{\xi}{4\pi G}\,\nabla^2\!\ln\lambda,
  \label{eq:RA_correction}
\end{equation}
where $\rho_A$ is the accumulated causal depth (recovering standard
GR in the equilibrium limit) and $\rho_\lambda$ is the gradient
correction term~\citep{Sandeman2026RADM}.

These departures from GR are not quantum-gravitational Planck-scale
effects. They are \emph{classical actualization-density gradient
effects} at cosmological and galactic scales:
\begin{itemize}[noitemsep]
  \item \textbf{Galactic rotation curves}: two complementary
    mechanisms give flat curves at all radii. In the disk, the 2D
    cylindrical Laplacian of $\ln\lambda$ gives $\nabla^2_{\rm 2D}\ln\lambda
    = -1/(R\,r_d)$, producing a $\Sigma_\lambda \propto 1/R$ surface density
    that yields a perfectly flat rotation curve with
    $v_{\rm flat}^2 = \xi/(2r_d)$. In the sparse halo ($\mu < 1$), DJW
    dimensional reduction gives $\Phi \propto \ln r$, also yielding
    $v = $ const. Both are RA-native derivations~\citep{Sandeman2026RADM}.
  \item \textbf{Hubble tension}: $H_{\rm eff}(x) \propto 1/\lambda(x)$
    predicts $H_{\rm local} \approx 73.1$~km/s/Mpc from the KBC void
    ($\delta \approx -0.20$), with a falsifiable prediction
    $H_{\rm Eridanus} \approx 75.9$~km/s/Mpc for the Eridanus
    supervoid ($\delta \approx -0.30$)~\citep{Sandeman2026RAGC}.
\end{itemize}

The precise sense in which $\Pact T_{\mu\nu} \to T_{\mu\nu}$ in
dense regions is: when $\lambda_{\rm local} \gg \lambda_c$ (the
critical actualization density at which the causal graph becomes
locally 3-dimensional), the off-shell fraction $1 - \Pact$ is
negligible and the standard GR source is recovered. The critical
density $\lambda_c$ is identified with the $\mu = 1$
Erd\H{o}s-R\'{e}nyi giant-component threshold in the Poisson-CSG
dynamics, connecting the cosmological and galactic departures from GR
to the same underlying graph-transition
phenomenon~\citep{Sandeman2026RAQI}.

% ── 7. Lean Verification Status ──────────────────────────────────────────────
\section{Lean~4 Verification Status}
\label{sec:lean}

The ingredients of Theorem~\ref{thm:main} have the following Lean~4
verification status:

\begin{center}
\renewcommand{\arraystretch}{1.25}
\begin{tabular}{lll}
\toprule
Ingredient & File & Status \\
\midrule
LLC (Definition~\ref{def:LLC}) & \texttt{RA\_Proofs\_Lean4.lean} & Lean-proved, 0 sorry \\
Graph Cut Theorem & \texttt{RA\_Proofs\_Lean4.lean} & Lean-proved, 0 sorry \\
Frame independence of $\dS$ & \texttt{RA\_Proofs\_Lean4.lean} & Lean-proved, 0 sorry \\
Vacuum suppression~\eqref{eq:vacuum_suppression} & (proved in RAGC) & Theorem \\
$\Pact$ conservation (Theorem~\ref{thm:P_act}) & \texttt{RA\_D1\_Proofs.lean} §16 & Structural proof \\
BDG locality (Lemma~\ref{lem:locality}) & \texttt{RA\_BDG\_locality.py} & Numerically verified \\
Rideout-Sorkin density limit & \citep{RideoutSorkin2000} & Published theorem \\
Benincasa-Dowker & \citep{Benincasa2010} & Published theorem \\
Lovelock uniqueness & \citep{Lovelock1971} & Published theorem \\
Bianchi identity & Differential geometry & Geometric identity \\
\bottomrule
\end{tabular}
\end{center}

\noindent
The $\Pact$ conservation proof (Section~\ref{sec:P_act}) is a
complete mathematical argument. Its Lean~4 formalisation in
\texttt{RA\_D1\_Proofs.lean}~§16 captures the logical structure with
structural theorems whose proofs reduce to the imported Lean results
(LLC, frame independence) and standard distributional analysis. Full
formalisation of the covariant distributional derivative step awaits
extension of the Lean~4 spacetime manifold library to distributional
tensor fields. This gap is narrower than the \texttt{amplitude\_locality}
axiom in the existing RA Lean files and is standard for mathematical
physics at this level.

The combined RA Lean~4 corpus now contains \textbf{101 results}
across \texttt{RA\_Proofs\_Lean4.lean} (26 results) and
\texttt{RA\_D1\_Proofs.lean} (75 results), with zero \texttt{sorry}
tags in either file.

% ── 8. Related Work ───────────────────────────────────────────────────────────
\section{Relation to Prior Work}
\label{sec:related}

\paragraph{Causal set theory and the BDG action.}
The Benincasa-Dowker action~\citep{Benincasa2010} was introduced as a
discrete scalar curvature for causal sets. Its path-integral
properties have been studied by Dowker, Surya, and
collaborators~\citep{Dowker2013}. The standard derivation of GR from
causal set theory proceeds by varying the BDG action, which yields
the Einstein equations in the continuum limit. The present derivation
differs in two respects: it uses the LLC as the starting point (a
vertex-level balance condition rather than an action principle), and
it arrives at a stronger conclusion (uniqueness of the field equation
with $\Lambda = 0$) by combining the BDG action with Lovelock's theorem
and the vacuum suppression result.

\paragraph{Jacobson's thermodynamic derivation.}
Jacobson~\citep{Jacobson1995} derived the Einstein equations from
thermodynamic considerations applied to local Rindler causal horizons.
That derivation also obtains the EFEs without assuming them, but
uses thermodynamic entropy as the primitive rather than a discrete
causal graph. The RA derivation uses the LLC as the discrete analogue
of the entropy condition; the RAGC Jacobson coefficient
$\alpha = 2\pi$ is confirmed to within 2\% using the full baryonic
disk mass~\citep{Sandeman2026RAGC,BlandHawthorn2016}.

\paragraph{Lovelock and uniqueness.}
The application of Lovelock's theorem to eliminate free parameters in
derived field equations has been noted before (see e.g.~\citep{Padmanabhan2015}).
The novelty here is the combination with the $\Pact$ conservation
theorem and the vacuum suppression result to achieve $\Lambda = 0$
without fine-tuning.

\paragraph{The cosmological constant.}
The cosmological constant problem --- why the observed $\Lambda$ is
$\sim10^{120}$ smaller than QFT vacuum energy estimates --- is
dissolved in RA by the vacuum energy suppression theorem: the vacuum
simply does not gravitate, because the metric is sourced only by
on-shell actualized states. This is a structural consequence of the
actualization ontology, not a fine-tuning. A related mechanism appears
in Padmanabhan's emergent gravity programme~\citep{Padmanabhan2015},
though the starting points differ substantially.

% ── 9. Conclusion ─────────────────────────────────────────────────────────────
\section{Conclusion}
\label{sec:conclusion}

We have proved that the Einstein field equations with $\Lambda = 0$
are the unique macroscopic field equation of the Relational Actualism
causal graph. The proof chain is:
\begin{enumerate}[noitemsep]
  \item LLC (Lean-verified) $\;\Rightarrow\;$
        $\nabla_\mu(\Pact T^{\mu\nu}) = 0$
        (Theorem~\ref{thm:P_act}, proved here)
  \item BDG locality: $\tau_{\rm BDG} \propto \lambda^{-1/4} \to 0$
        (Lemma~\ref{lem:locality}, proved here)
  \item $\langle\SBDG\rangle \to (l_P^2/4)R$
        (Benincasa-Dowker~\citep{Benincasa2010})
  \item $H_{\mu\nu} = \alpha G_{\mu\nu} + \Lambda g_{\mu\nu}$
        (Lovelock~\citep{Lovelock1971})
  \item $\Lambda = 0$ (vacuum suppression~\citep{Sandeman2026RAGC})
  \item $\nabla_\mu G^{\mu\nu} \equiv 0$ (Bianchi, geometric identity)
\end{enumerate}
General relativity is an output of RA, not an input. It is the unique
consistent description of the RA-CSG dynamics in the regime where the
actualization projector $\Pact$ acts as the identity on the stress
tensor. The departures from GR in sparse regions --- galactic rotation
curves and the Hubble tension --- are predicted by the same framework
that derives GR, as leading-order corrections from the spatial
gradient of the actualization density field.

\section*{Acknowledgements}

The author is an independent researcher with no institutional affiliation.
This research received no external funding.
The author declares no conflicts of interest.

\smallskip\noindent
\textbf{AI use disclosure.}
Claude (Anthropic, model \texttt{claude-sonnet-4-6}) was used in the
preparation of this manuscript, including drafting and revision of text,
mathematical derivation assistance and verification scaffolding,
\LaTeX{} compilation, and Lean~4 proof development.
The author is solely responsible for all scientific content,
mathematical claims, and conclusions.
AI tools were used in accordance with IOP Publishing's generative AI policy.
Gemini (Google DeepMind, Gemini~2.0~Pro) was used for supplementary
literature checking and cross-verification of symbolic calculations.

\smallskip\noindent
\textbf{Data availability.}
All Lean~4 proof files and Python verification scripts supporting this
article are openly available at
\url{https://github.com/jsandeman/Relational-Actualism}.

\smallskip\noindent
\textbf{Preprint.}
A preprint version of this article is available at
\href{https://doi.org/10.5281/zenodo.19270020}{doi:10.5281/zenodo.19270020}.

\begin{thebibliography}{99}

\bibitem[Benincasa \& Dowker(2010)]{Benincasa2010}
D.~M.~T. Benincasa and F.~Dowker.
\newblock The scalar curvature of a causal set.
\newblock \textit{Phys.\ Rev.\ Lett.}, 104:181301, 2010.
\newblock \href{https://doi.org/10.1103/PhysRevLett.104.181301}{doi:10.1103/PhysRevLett.104.181301}.

\bibitem[Bland-Hawthorn \& Gerhard(2016)]{BlandHawthorn2016}
J.~Bland-Hawthorn and O.~Gerhard.
\newblock The galaxy in context: structural, kinematic, and integrated properties.
\newblock \textit{Annu.\ Rev.\ Astron.\ Astrophys.}, 54:529--596, 2016.

\bibitem[Dowker(2013)]{Dowker2013}
F.~Dowker.
\newblock Introduction to causal sets and their phenomenology.
\newblock \textit{Gen.\ Relativ.\ Gravit.}, 45:1651--1667, 2013.

\bibitem[Jacobson(1995)]{Jacobson1995}
T.~Jacobson.
\newblock Thermodynamics of spacetime: the Einstein equation of state.
\newblock \textit{Phys.\ Rev.\ Lett.}, 75:1260--1263, 1995.

\bibitem[Lovelock(1971)]{Lovelock1971}
D.~Lovelock.
\newblock The Einstein tensor and its generalizations.
\newblock \textit{J.\ Math.\ Phys.}, 12:498--501, 1971.

\bibitem[Meiburg et al.(2025)]{Meiburg2025LQI}
J.~Meiburg, E.~Lessa, and G.~Soldati.
\newblock Lean-QuantumInfo: a library for quantum information in Lean~4.
\newblock \textit{arXiv:2510.08672}, 2025.

\bibitem[Padmanabhan(2015)]{Padmanabhan2015}
T.~Padmanabhan.
\newblock Emergent gravity paradigm: recent progress.
\newblock \textit{Mod.\ Phys.\ Lett.~A}, 30:1540007, 2015.

\bibitem[Rideout \& Sorkin(2000)]{RideoutSorkin2000}
D.~P. Rideout and R.~D. Sorkin.
\newblock A classical sequential growth dynamics for causal sets.
\newblock \textit{Phys.\ Rev.~D}, 61:024002, 2000.

\bibitem[Sandeman(2026RAQM)]{Sandeman2026RAQM}
J.~F. Sandeman.
\newblock {Relational Actualism and Quantum Mechanics (RAQM)}.
\newblock \textit{Foundations of Physics} (under review), 2026.
\newblock \href{https://doi.org/10.5281/zenodo.19174380}{doi:10.5281/zenodo.19174380}.

\bibitem[Sandeman(2026RAGC)]{Sandeman2026RAGC}
J.~F. Sandeman.
\newblock {Relational Actualism: Gravity and Cosmology (RAGC)}.
\newblock \textit{Preprint}, 2026.
\newblock \href{https://doi.org/10.5281/zenodo.19197999}{doi:10.5281/zenodo.19197999}.

\bibitem[Sandeman(2026RAEB)]{Sandeman2026RAEB}
J.~F. Sandeman.
\newblock {Relational Actualism: The Engine of Becoming (RAEB)}.
\newblock \textit{Preprint}, 2026.
\newblock \href{https://doi.org/10.5281/zenodo.19198120}{doi:10.5281/zenodo.19198120}.

\bibitem[Sandeman(2026RASM)]{Sandeman2026RASM}
J.~F. Sandeman.
\newblock {Relational Actualism and the Standard Model (RASM)}.
\newblock \textit{Preprint}, 2026.
\newblock \href{https://doi.org/10.5281/zenodo.19229335}{doi:10.5281/zenodo.19229335}.

\bibitem[Sandeman(2026RATM)]{Sandeman2026RATM}
J.~F. Sandeman.
\newblock {Relational Actualism: The Topology of Matter (RATM)}.
\newblock \textit{Preprint}, 2026.
\newblock \href{https://doi.org/10.5281/zenodo.19265651}{doi:10.5281/zenodo.19265651}.

\bibitem[Sandeman(2026RADM)]{Sandeman2026RADM}
J.~F. Sandeman.
\newblock {Relational Actualism: Dark Matter as Actualization-Density Topology (RADM)}.
\newblock \textit{Preprint}, 2026.
\newblock \href{https://doi.org/10.5281/zenodo.19198513}{doi:10.5281/zenodo.19198513}.

\bibitem[Sandeman(2026RAQI)]{Sandeman2026RAQI}
J.~F. Sandeman.
\newblock {Relational Actualism: Implications for Quantum Information (RAQI)}.
\newblock \textit{Preprint}, 2026.
\newblock \href{https://doi.org/10.5281/zenodo.19198285}{doi:10.5281/zenodo.19198285}.

\end{thebibliography}

\end{document}
